\documentclass[10pt]{article}
\usepackage[a4paper,margin=20mm]{geometry}
\usepackage{amsmath,amssymb}
\usepackage[T1]{fontenc}
\usepackage{hyperref}
\title{A cubic-box counterexample to Conjecture 425}
\author{Local AI-assisted solution draft}
\date{3 October 2026}
% Compact consistent title for a short mathematical note.
\makeatletter
\renewcommand{\maketitle}{\begin{center}
{\Large\bfseries\@title\par}\vspace{0.6em}
{\small\@author\par\@date}
\end{center}\vspace{0.5em}}
\makeatother
\setlength{\emergencystretch}{3em}
\usepackage{url}
\begin{document}
\maketitle
\paragraph{Source and interpretation.}
Conjecture \texttt{00000000425} (original statement and pinned commit in
\texttt{SOURCE.md}) asserts that the
coefficient sequence in the diagonal specialization of the boxed plane
partition generating function is always log-concave. We use the standard
volume generating polynomial, with the diagonal specialization meaning equal
box side lengths. The counterexample is the $2\times2\times2$ box.

\paragraph{Definitions.}
A plane partition in this box is an array
$\left(\begin{smallmatrix}a&b\\c&d\end{smallmatrix}\right)$ with entries in
$\{0,1,2\}$, weakly decreasing along both rows and columns:
$a\ge b$, $a\ge c$, $b\ge d$, and $c\ge d$.
Its volume is $a+b+c+d$. Let $c_k$ count such arrays of volume $k$.
Log-concavity requires $c_k^2\ge c_{k-1}c_{k+1}$ at every interior index.

\paragraph{Counterexample and proof.}
There is exactly one array of volume zero, namely the zero array, and exactly
one of volume one, with a single $1$ in the top-left entry. The three arrays
of volume two are
\[
\begin{pmatrix}2&0\\0&0\end{pmatrix},\qquad
\begin{pmatrix}1&1\\0&0\end{pmatrix},\qquad
\begin{pmatrix}1&0\\1&0\end{pmatrix}.
\]
These are exhaustive: either one entry equals $2$, which by monotonicity
must be the top-left entry, or two entries equal $1$, giving exactly the
two displayed order ideals. Consequently
\[
c_0=1,\qquad c_1=1,\qquad c_2=3,\qquad
c_1^2=1<3=c_0c_2.
\]
Thus the asserted universal log-concavity is false. It is unnecessary to
consider the additional equality characterization after this first clause
has failed.

\paragraph{Independent polynomial check.}
The complete polynomial, obtained both by direct enumeration and by the
MacMahon product, is
\[
\prod_{i,j,k=1}^{2}\frac{1-q^{i+j+k-1}}{1-q^{i+j+k-2}}
=1+q+3q^2+3q^3+4q^4+3q^5+3q^6+q^7+q^8.
\]
The accompanying Python verifier independently computes the product as an
exact integer formal power series and enumerates the arrays.

\paragraph{Direct formal connection to the MacMahon product.}
Write its numerator and denominator as $N(q)$ and $D(q)$. For the cubic box
of side two, the exact triple products give
\[
 D(q)=1-q-3q^2+O(q^3),\qquad N(q)=1-q^2+O(q^3).
\]
Since $D(0)=1$, formal division defines a coefficient sequence $f$ with
$D(q)F(q)=N(q)$. Comparing the first three coefficients gives
$f_0=1$, $f_1-f_0=0$, and $f_2-f_1-3f_0=-1$; thus again
$(f_0,f_1,f_2)=(1,1,3)$. This directly refutes log-concavity for the product
specified in the conjecture, without using the counting formula as an
unproved assumption.

The self-contained Lean 4.19.0 file \texttt{Main.lean} formalizes both routes.
It verifies the complete, duplicate-free enumeration of the 81 arrays and
the precise plane-partition predicate. It also defines integer formal series,
Cauchy multiplication, and the actual three nested MacMahon products. The
terminating recursion \texttt{formalQuotient} constructs every coefficient;
\texttt{formalQuotient\_spec} proves $D F=N$ at every index whenever $D(0)=1$.
The theorem \texttt{macmahon\_product\_counterexample} therefore refutes
log-concavity for an explicitly constructed infinite coefficient sequence,
not for a predicate whose existence is merely assumed. Its first three
coefficients are also proved equal to the direct enumeration above.
All calculations use kernel proofs, with no placeholders or custom axioms.

\paragraph{Alternative reading of ``coefficientwise.''}
If the intended claim instead compares consecutive polynomials $F_n(q)$
for $n\times n\times n$ boxes, it is false as well. The completely filled
box is the unique term of maximum volume, so $F_n$ is monic of degree
$n^3$. At $n=2$, the polynomial $F_2(q)^2-F_1(q)F_3(q)$ has coefficient
$-1$ at $q^{28}$, because the two product degrees are $16$ and $28$.
This additional argument covers that reading in ordinary mathematics;
the Lean artifact formalizes the scalar coefficient-sequence reading above.
The source does not explicitly define ``diagonal specialization,'' so we
retain this interpretation boundary rather than claiming all possible
meanings have been formalized.

\end{document}
